% defer/whichtochoose.tex
% mainfile: ../perfbook.tex
% SPDX-License-Identifier: CC-BY-SA-3.0

\section{如何选择？}
\label{sec:defer:Which to Choose?}
%
\epigraph{总是选择看起来最好的道路，无论它多么崎岖；
	  习惯很快会使它变得轻松而愉快。}
	  {Pythagoras}

\Cref{sec:defer:Which to Choose? (Overview)}
提供了一个高层概述，然后
\cref{sec:defer:Which to Choose? (Details)}
提供了本章中介绍的延迟处理技术之间差异的更详细视图。
本讨论假设一个足够大的链式数据结构，读者不会在一次遍历到
另一次遍历之间持有引用，并且元素可能在任何位置和任何时间
被添加到结构中或从中移除。
\Cref{sec:defer:Which to Choose? (Production Use)}
然后指出了\IXpl{hazard pointer}、sequence locking和RCU的
一些公开可见的生产用途\@。
本讨论应该帮助你在这些技术之间做出明智的选择。

\subsection{如何选择？（概述）}
\label{sec:defer:Which to Choose? (Overview)}

\begin{table*}
\rowcolors{1}{}{lightgray}
\renewcommand*{\arraystretch}{1.25}
\footnotesize
\centering\OneColumnHSpace{-.3in}
\ebresizewidth{
\begin{tabularx}{5.3in}{>{\raggedright\arraybackslash}p{1.1in}
    >{\raggedright\arraybackslash}p{1.0in}
    >{\raggedright\arraybackslash}X
    >{\raggedright\arraybackslash}X
    >{\raggedright\arraybackslash}p{.9in}}
	\toprule
	属性
		& Reference Counting
			& Hazard Pointers
				& Sequence Locks
					& RCU \\
%		  RC	  HP	  SL	  RCU \\
	\midrule
	读者
		& 慢且不可扩展
			& 快且可扩展
				& 快且可扩展
					& 快且可扩展 \\
	内存开销
		& 每个对象一个计数器
			& 每个读者每个对象一个指针
				& 无保护
					& 无 \\
	保护持续时间
		& 可以很长
			& 可以很长
				& 无保护
					& 用户必须限制持续时间 \\
	遍历重试的需要
		& 如果对象被删除
			& 如果对象被删除
				& 如果有任何更新
					& 从不 \\
	回收时机
		& 立即
			& 批处理延迟
				& 不适用
					& 已有读者完成加上批处理延迟 \\
	\bottomrule
\end{tabularx}
}
\caption{选择哪种延迟处理技术？（概述）}
\label{tab:defer:Which Deferred Technique to Choose? (Overview)}
\end{table*}

\Cref{tab:defer:Which Deferred Technique to Choose? (Overview)}
显示了一些区分延迟回收技术的高层属性。

"读者"行总结了
\cref{fig:defer:Pre-BSD Routing Table Protected by RCU QSBR}
中呈现的结果，表明除了\IXalt{reference counting}{reference count}
之外，其他技术都享有相当快速和可扩展的读者。

"内存开销"行评估了每种技术对外部存储的需求，
用于记录读者保护信息。
RCU依赖quiescent state，因此不需要存储来表示读者，
无论是在对象内部还是外部。
Reference counting可以在结构中每个对象内使用一个整数，
不需要额外存储。
Hazard pointer需要配置对象外部的指针，并且必须有足够的指针
供每个CPU或线程在任何给定时间跟踪所有被引用的对象。
鉴于大多数基于hazard pointer的遍历只需要少量hazard pointer，
这在实践中通常不是问题。
当然，sequence lock不提供指针遍历保护，
这就是它们通常用于静态数据的原因。

\QuickQuiz{
	为什么用户不能在需要时动态分配hazard pointer？
}\QuickQuizAnswer{
	他们可以，但代价是额外的读者遍历开销，
	并且在某些环境中需要处理内存分配失败。
}\QuickQuizEnd

"保护持续时间"行描述了用户可以保护给定对象的时间长度的
约束（如果有的话）。
Reference counting和hazard pointer都可以长时间保护对象
而不会产生不良副作用，但对即使一个对象维持RCU引用也会
阻止所有其他RCU保护的对象被释放。
因此，RCU读者必须相对较短以避免耗尽系统内存，
SRCU、Tasks RCU和Tasks Trace RCU等专用实现是此规则的例外。
同样，sequence lock不提供指针遍历保护，
这就是它们通常用于静态数据的原因。

"遍历重试的需要"行说明了是否可以无条件地获取对给定对象的
新引用（如RCU那样），还是引用获取可能失败从而导致重试操作
（如reference counting、hazard pointer和sequence lock的情况）。
在reference counting和hazard pointer的情况下，
只有在尝试获取对给定对象的引用时该对象正在被删除过程中失败，
才需要重试，这个主题在下一节中有更详细的讨论。
Sequence locking当然必须在其critical section与任何更新并发
运行时重试该critical section。

\QuickQuiz{
	但是Linux内核的\co{kref}引用计数器不是允许
	保证无条件引用获取吗？
}\QuickQuizAnswer{
	是的，它们允许，但这个保证只在已经持有引用的情况下
	无条件适用。
	考虑到这一点，请回顾
	\cref{sec:defer:Which to Choose?}开头的段落，
	特别是"足够大以至于读者不会在一次遍历到另一次遍历之间
	持有引用"的部分。
}\QuickQuizEnd

"回收时机"行给出了从最后一个读者使用完被移除对象到
该被移除对象可以被释放之间的最小延迟。
Reference counting是唯一能够在最后一个读者完成后立即释放的
技术。
这一优势当然是reference counting大缺点的另一面——
它的读者慢且不可扩展。
理论上，hazard pointer也可以立即回收，但在实践中
生产级hazard pointer实现使用批处理来将扫描hazard pointer的
开销分摊到多次更新上，这种批处理会产生额外延迟。
RCU必须等待所有已有读者完成，无论这些读者是否与
新移除的对象有任何关系，然后，与hazard pointer一样，
生产级RCU实现由于批处理而产生额外延迟。

当然，不同的行在不同情况下会有不同的重要性级别。
例如，如果你当前的代码在使用hazard pointer时遇到读侧
可扩展性问题，那么hazard pointer可能需要重试引用获取
这一点就无关紧要了，因为你当前的代码已经处理了这种情况。
类似地，如果响应时间的考虑已经限制了读者遍历的持续时间
（在内核和底层应用中经常如此），那么RCU有持续时间限制要求
就无关紧要了，因为你的代码已经满足了这些要求。
同样，如果读者必须已经写入它们正在遍历的对象，
引用计数器的读侧开销可能就不那么重要了。
当然，如果要保护的数据在静态分配的变量中，
那么sequence locking无法保护指针这一点就无关紧要了。

最后，有一些关于根据延迟动态采样在hazard pointer和RCU之间
动态切换的工作~\cite{Balmau:2016:FRM:2935764.2935790}。
这将hazard pointer和RCU之间的选择推迟到运行时，
并将决策责任委托给软件。

尽管如此，这张表在选择这些技术时应该非常有帮助。
但那些希望了解更多细节的读者应该继续阅读下一节。

\subsection{如何选择？（详情）}
\label{sec:defer:Which to Choose? (Details)}

\begin{table*}
\rowcolors{1}{}{lightgray}
\renewcommand*{\arraystretch}{1.25}
\footnotesize
\centering\OneColumnHSpace{-.3in}
\ebresizewidth{
\begin{tabularx}{5.3in}{>{\raggedright\arraybackslash}p{1.1in}
    >{\raggedright\arraybackslash}p{1.2in}
    >{\raggedright\arraybackslash}X
    >{\raggedright\arraybackslash}X
    >{\raggedright\arraybackslash}p{.9in}}
	\toprule
	属性
		& Reference Counting
			& Hazard Pointers
				& Sequence Locks
					& RCU \\
%		  RC	  HP	  SL	  RCU \\
	\midrule
	存在性保证
		& 复杂
			& 是
				& 否
					& 是 \\
	更新与读者并发推进
		& 是
			& 是
				& 否
					& 是 \\
	读者间争用
		& 高
			& 无
				& 无
					& 无 \\
	读者每个critical section开销
		& 不适用
			& 不适用
				& 两次\tco{smp_mb()}
					& 从无到两次\tco{smp_mb()} \\
	读者每个对象遍历开销
		& 读-修改-写原子操作、memory\-/barrier指令和缓存未命中
			& \tco{smp_mb()}\parnote[*]{通过使用Linux内核的
				\tco{membarrier()}系统调用，此\tco{smp_mb()}
				可以降级为编译器\tco{barrier()}。}
				& 无，但不安全
					& 无（volatile访问） \\
	读者前进保证
		& Lock free
			& Lock free
				& 阻塞
					& Bounded wait free（QSBR \& urcu） \\
	读者引用获取
		& 可能失败（有条件）
			& 可能失败（有条件）
				& 不安全
					& 不会失败（无条件） \\
	内存占用
		& 有界
			& 有界
				& 有界
					& 无界 \\
	回收前进保证
		& Lock free
			& Lock free
				& 不适用
					& 阻塞 \\
	自动回收
		& 是
			& 视用例
				& 不适用
					& 视用例 \\
	代码行数
		& 94
			& 79
				& 79
					& 73 \\
	\bottomrule
\end{tabularx}
}
\begin{minipage}{\onecolumntextwidth}
	\vspace*{1ex}
	\parnotes
\end{minipage}
\caption{选择哪种延迟处理技术？（详情）}
\label{tab:defer:Which Deferred Technique to Choose?  (Details)}
\end{table*}

\Cref{tab:defer:Which Deferred Technique to Choose? (Details)}
提供了更详细的经验法则，可以帮助你在本章介绍的四种延迟处理
技术之间做出选择。

如"存在性保证"行所示，
如果你需要链式数据元素的\IXpl{existence guarantee}，
你必须使用reference counting、hazard pointer或RCU\@。
Sequence lock不提供存在性保证，而是提供更新检测，
重试任何遇到更新的read-side critical section。

当然，如"更新与读者并发推进"行所示，
这种更新检测意味着sequence locking不允许更新者和读者
并发地向前推进。
毕竟，阻止这种并发前进正是使用sequence locking的全部意义！
这种情况指出了一种方式：将sequence locking与reference counting、
hazard pointer或RCU结合使用，以同时提供存在性保证和更新检测。
事实上，Linux内核在路径名查找期间就以这种方式结合了RCU和
sequence locking。

"读者间争用"、"读者每个critical section开销"和
"读者每个对象遍历开销"行给出了这些技术读侧开销的大致感觉。
Reference counting的开销可能相当大，读者之间存在争用，
而且每遍历一个对象都需要一个完全有序的读-修改-写原子操作。
Hazard pointer为每个遍历的数据元素产生一次\IX{memory barrier}的开销，
sequence lock为每次尝试执行critical section产生一对
memory barrier的开销。
RCU实现的开销从零到每个read-side critical section一对
memory barrier不等，因此RCU提供了最佳性能，
特别是对于遍历许多数据元素的read-side critical section。
当然，所有延迟处理变体的读侧开销都可以通过批处理来减少，
使得每个读侧操作覆盖更多数据。

\QuickQuiz{
	但是\cref{sec:defer:Hazard Pointers}中某个小测验的
	答案不是说非对称屏障对可以消除hazard pointer的
	读侧\co{smp_mb()}吗？
}\QuickQuizAnswer{
	是的，确实如此。
	然而，这样做可以被认为将hazard pointer的
	"回收前进保证"行（稍后讨论）从lock-free变为阻塞，
	因为在内核中禁用中断自旋的CPU将阻止非对称屏障的
	更新侧部分完成。
	在Linux内核中，理论上可以通过以下方式防止此类阻塞：
	使用\co{CONFIG_NO_HZ_FULL}构建内核，
	在启动时将相关CPU指定为\co{nohz_full}，
	确保给定CPU上同一时间只有一个线程可运行，
	并避免调用内核。
	或者，你可以确保内核没有任何可能导致CPU在禁用中断
	状态下自旋的bug。

	鉴于在Linux内核中CPU禁用中断自旋似乎相当罕见，
	人们可以反驳说非对称屏障hazard pointer更新在实践中
	是非阻塞的，即使在理论上不是。
}\QuickQuizEnd

"读者前进保证"行显示只有RCU具有
\IXalth{bounded wait-free}{bounded}{wait free}
\IX{forward-progress guarantee}，这意味着
它可以通过执行有界数量的指令来完成有限遍历。

\QuickQuiz{
	等等！！！
	Linux内核RCU的\co{rcu_read_unlock()}可能进行锁获取，
	这不可能是wait-free的！
	你在搞什么？？？
}\QuickQuizAnswer{
	你似乎错过了那个单元格中的"（QSBR \& urcu）"。

	非抢占式Linux内核使用QSBR，其中\co{rcu_read_lock()}
	不仅是bounded population-oblivious wait free，
	而且执行零条指令。

	用户空间RCU执行有限数量的简单指令序列，
	因此也符合bounded population-oblivious wait free。

	一个例外是用户空间RCU的"防弹"变体，其中第一次调用
	\co{rcu_read_lock()}会将当前线程注册为RCU读者。
	但只有第一次调用如此，因此可以通过确保每个线程在
	执行任何需要wait-free前进保证的工作之前
	调用\co{rcu_read_lock()}来消除这一额外开销。

	但是可抢占式Linux内核呢？

	在那里，目标是激进的实时响应，这在
	\cref{sec:advsync:Parallel Real-Time Computing}中讨论。

	与此同时，请注意Linux内核只支持有限数量的CPU，
	并且\co{rcu_preempt_deferred_qs_irqrestore()}中的
	锁获取（从\co{rcu_read_unlock()}间接调用）
	是在禁用中断的情况下进行的。
	此外，这些锁是复制的，因此给定的锁通常只被十六个CPU
	中的一个获取。
	这限制了锁争用，并防止抢占和中断在持有这些锁时
	延迟执行，从而提供了出色的前进
	保证~\cite{BjoernBrandenburgPhD,DanielBristot2019RTtrace,Hillier86,DipankarSarma2004OLSscalability}。

	简而言之，wait-free同步等形式化概念有其用武之地，
	但真正重要的是弄清楚你的用户真正需要什么，
	然后为他们实现它。
}\QuickQuizEnd

"读者引用获取"行表明只有RCU能够无条件获取引用。
Sequence lock的条目是"不安全"，因为sequence lock检测更新
而不是获取引用。
Reference counting和hazard pointer都要求在给定获取失败时
从头重新开始遍历。
要理解这一点，考虑一个包含对象~A、B、C和~D（按此顺序）的
链表，以及以下一系列事件：

\begin{enumerate}
\item	读者获取对对象~B的引用。
\item	更新者移除对象~B，但因为读者持有引用而不释放它。
	列表现在包含对象~A、C和~D，
	对象~B的\co{->next}指针被设置为\co{HAZPTR_POISON}。
\item	更新者移除对象~C，使得列表现在包含对象~A和~D\@。
	因为没有对象~C的引用，它被立即释放。
\item	读者试图前进到现已被移除的对象~B之后的对象的后继者，
	但被毒化的\co{->next}指针阻止了这一操作。
	这是一件好事，因为否则对象~B的\co{->next}指针将指向
	空闲列表。
\item	因此读者必须从列表头部重新开始遍历。
\end{enumerate}

因此，当引用获取失败时，hazard pointer或reference counter遍历
必须从头重新开始该遍历。
在嵌套链式数据结构的情况下，例如包含链表的树，
遍历必须从最外层数据结构重新开始。
这种情况给予了RCU显著的易用性优势。

然而，RCU的易用性优势并非免费的，
如"内存占用"行所示。
RCU支持无条件引用获取意味着它必须避免释放给定RCU读者
可到达的任何对象，直到该读者完成。
因此，RCU具有无界的内存占用，至少在更新没有被限流的情况下。
相比之下，reference counting和hazard pointer只需要保留
那些实际被并发读者引用的数据元素。

这种内存占用和获取失败之间的张力有时在Linux内核中
通过结合使用RCU和引用计数器来解决。
RCU用于短期引用，这意味着RCU read-side critical section可以很短。
这些短的RCU read-side critical section反过来意味着相应的
RCU \IXpl{grace period}也可以很短，从而限制了内存占用。
对于少数需要更长期引用的数据元素，使用reference counting。
这意味着引用获取失败的复杂性只需要处理那些少数数据元素：
大部分引用获取是无条件的，这要感谢RCU\@。
更多关于将reference counting与其他同步机制结合使用的信息
见\cref{sec:together:Refurbish Reference Counting}。

"回收前进保证"行显示hazard pointer可以提供
非阻塞更新~\cite{MagedMichael04a,HerlihyLM02}。
Reference counting可能可以也可能不可以，取决于实现。
然而，sequence locking由于其更新侧锁而无法提供非阻塞更新。
RCU更新者必须等待读者，这也排除了完全非阻塞的更新。
然而，在某些情况下唯一的阻塞操作是等待释放内存，
这导致了一种在很多目的上与非阻塞一样好的
情况~\cite{MathieuDesnoyers2012URCU}。

如"自动回收"行所示，只有reference counting可以自动释放内存，
而且仅限于非循环数据结构。
Hazard pointer和RCU的某些用例可以使用\emph{link count}提供
自动回收，link count可以被视为引用计数，
但仅适用于来自数据结构其他部分的入向链接~\cite{MagedMichael2018FollyHazptr}。

最后，"代码行数"行显示了Pre-BSD路由表实现的大小，
大致反映了相对易用性。
也就是说，重要的是要注意reference counting和sequence locking的
实现是有bug的，而正确的reference counting实现要复杂得多~\cite{Valois95a,MagedMichael95a}。
就sequence locking而言，正确的实现需要添加某种其他同步机制，
例如hazard pointer或RCU，以便sequence locking检测并发更新，
而另一种机制提供安全的引用获取。

随着使用这些技术（无论是单独使用还是组合使用）获得更多经验，
本节中列出的经验法则将需要改进。
然而，本节确实反映了当前的技术水平。

\subsection{如何选择？（生产使用）}
\label{sec:defer:Which to Choose? (Production Use)}

本节指出了hazard pointer、sequence locking和RCU的一些
公开可见的生产用途。
Reference counting被省略，不是因为它不重要，而是因为它不仅
被广泛使用，而且在回溯半个世纪的教科书中都有大量记载。
列出这些其他技术的生产用途的期望好处之一是提供可供研究
（或发现bug，视情况而定）的示例。\footnote{
	感谢Mathias Stearn、Matt Wilson、David Goldblatt、LiveJournal
	用户fanf、Nadav Har'El、Avi Kivity、Dmitry Vyukov、Raul Guitterez
	S.、Twitter用户@peo3、Paolo Bonzini和Thomas Monjalon
	找到了大量这些用例。}

\subsubsection{风险指针的生产用途}
\label{sec:defer:Production Uses of Hazard Pointers}

2010年，Keith Bostic向WiredTiger添加了hazard pointer的一个
变体~\cite{KeithBostic2010WiredTigerhazptr}。
2015年发布的MongoDB 3.0包含了WiredTiger，因此也包含了hazard pointer。

2011年，Samy Al Bahra向Concurrency Kit库添加了
hazard pointer~\cite{SamyAlBahra2011ckhp}。

2014年，Maxim Khizhinsky向libcds添加了
hazard pointer~\cite{MaximKhizhinsky2014libcdsHazptr}。

2015年，David Gwynne向OpenBSD引入了shared reference pointer，
一种hazard pointer的形式~\cite{DavidGwynne2015srp}。

2017--2018年，Rust语言的
\co{arc-swap}~\cite{MichalVaner2018arc-swapHazptr}和
\co{conc}~\cite{crates.io.user.ticki2017concHazptr}
crate各自实现了自己的hazard pointer。

2018年，Maged Michael向Facebook的Folly库添加了
hazard pointer~\cite{MagedMichael2018FollyHazptr}，
并在该库中被大量使用。

\subsubsection{顺序锁的生产用途}
\label{sec:defer:Production Uses of Sequence Locking}

Linux内核在2003年的v2.5.60中添加了
sequence locking~\cite{JonathanCorbet2003seqlock}，
这是从x86实现\co{gettimeofday()}系统调用时使用的一种
临时技术泛化而来。

2011年，Samy Al Bahra向Concurrency Kit库添加了
sequence locking~\cite{SamyAlBahra2011ckseqlock}。

2013年，Paolo Bonzini向QEMU模拟器添加了简单的
sequence-lock~\cite{PaoloBonzini2013QEMUseqlock}。

2016年，Alexis Menard在Chromium中抽象了一个
sequence-lock实现~\cite{AlexisMenard2016ChromiumSeqLock}。

2018年，一个简单的sequence locking实现被添加到
\co{jemalloc()}中~\cite{DavidGoldblatt2018seqlock}。
eigen库也有一个由类似于sequence locking的机制管理的
特殊用途队列。

\subsubsection{RCU的生产用途}
\label{sec:defer:Production Uses of RCU}

IBM的VM/XA在1980年代的某个时候采用了passive serialization，
一种类似于RCU的机制~\cite{Hennessy89}。

DYNIX/ptx在1993年采用了RCU~\cite{McKenney98,Slingwine95}。

Linux内核在2002年采用了Dipankar Sarma的RCU
实现~\cite{Torvalds2.5.43}。

用户空间RCU项目于2009年启动~\cite{MathieuDesnoyers2009URCU}。

Knot DNS项目于2010年开始使用用户空间RCU
库~\cite{LubosSlovak2010KnotDNSRCU}。
同年，OSv内核添加了RCU
实现~\cite{AviKivity2013OSvRCU}，
后来添加了RCU保护的链表~\cite{AviKivity2013OSvRCUlist}
和RCU保护的哈希表~\cite{AviKivity2013OSvRCUhash}。

2011年，Samy Al Bahra向Concurrency Kit库添加了epoch
（RCU的一种形式~\cite{UCAM-CL-TR-579,KeirFraser2007withoutLocks}）~\cite{SamyAlBahra2011ckEpoch}。

NetBSD在2012年的v6.0中开始使用上述passive
serialization~\cite{NetBSD2012pserialize}。
其中，passive serialization被用于NetBSD数据包过滤器
（NPF）~\cite{MindaugasRasiukevicius2014NPFRCU}。

2015年，Paolo Bonzini通过用户空间RCU库的友好分支向QEMU模拟器
添加了RCU支持~\cite{MikeDay2013RCUqemu,PaoloBonzini2013QEMURCU}。

2015年，Maxim Khizhinsky向libcds添加了
RCU~\cite{MaxKhiszinsky2015C++RCU}。

Mindaugas Rasiukevicius于2016年实现了libqsbr，
其中包含\IXacr{qsbr}和\IXacrf{ebr}~\cite{MindaugasRasiukevicius2016libqsbr}，
两者都是RCU的实现类型。

Sheth等人~\cite{HarshalSheth2016goRCU}
展示了利用Go的垃圾收集器提供类RCU功能的价值，
Go编程语言提供了一种\co{Value}类型可以提供此
功能。\footnote{
	见\url{https://golang.org/pkg/sync/atomic/\#Value}，
	特别是"Example (ReadMostly)"。}

Matt Klein描述了Envoy Proxy中使用的一种类RCU
机制~\cite{MattKlein2017EnvoyRCU}。

Honnappa Nagarahalli于2018年向数据平面开发套件（DPDK）添加了
RCU库~\cite{HonnappaNagarahalli2018dpdkRCU}。

Stjepan Glavina将基于epoch的RCU实现合并到了Rust语言的
\co{crossbeam}并发支持"crate"集中~\cite{StjepanGlavina2018RustRCU}。

Jason Donenfeld在将WireGuard移植到Windows~NT内核的过程中
实现了一个RCU实现~\cite{JasonDonenfeld2021:WindowsNTwireguardRCU}。

最后，任何支持垃圾收集的并发语言（不仅仅是Go！\@）都可以
以零增量成本获得RCU实现的更新侧。

\subsubsection{生产用途总结}
\label{sec:defer:Summary of Production Uses}

也许有一天，sequence locking、hazard pointer和RCU都会像
引用计数器一样被广泛使用和熟知。
在那一天到来之前，这些机制的当前生产用途应该有助于指导
机制的选择，并展示如何最好地应用每种机制。
至此，我们揭开了\cpageref{sec:defer:Mysteries Which to choose}
上提出的最后一个谜题。

下一节讨论更新操作，这对本章描述的许多读多写少机制来说是
一个棘手的问题。
